\eventsection{CONSIDERA\c{C}\~OES FINAIS}

O estudo comparou algoritmos de \textit{bagging} e \textit{boosting} para detec\c{c}\~ao de surtos de mal\'aria no estado do Amazonas com base em um protocolo temporal rigoroso. Os resultados indicaram desempenho superior do LightGBM em rela\c{c}\~ao ao XGBoost e ao Random Forest, tanto em capacidade discriminativa quanto em qualidade probabil\'istica. Esse achado sugere que modelos de \textit{boosting} s\~ao mais adequados para apoiar sistemas de alerta em contextos epidemiol\'ogicos com baixa preval\^encia de surtos.

Como implica\c{c}\~ao pr\'atica, a aplica\c{c}\~ao desses modelos pode contribuir para orientar monitoramento preventivo, prioriza\c{c}\~ao territorial e revis\~ao de limiares de risco em rotinas de vigil\^ancia. No entanto, o uso operacional deve permanecer vinculado \`a avalia\c{c}\~ao t\'ecnica das equipes de sa\'ude, com monitoramento cont\'inuo de desempenho, recalibra\c{c}\~ao peri\'odica e transpar\^encia sobre os crit\'erios de decis\~ao.

Como trabalhos futuros, recomenda-se ampliar a valida\c{c}\~ao para diferentes recortes geogr\'aficos, explorar calibra\c{c}\~ao probabil\'istica formal e comparar o ajuste por peso de classe com t\'ecnicas de reamostragem e abordagens custo-sens\'iveis, como SMOTE e aprendizado custo-sens\'ivel \cite{chawla2002smote,elkan2001foundations}. Tamb\'em s\~ao relevantes estudos com m\'ultiplos anos, deriva de atributos, monitoramento de preval\^encia e estabilidade da taxa de alertas por munic\'ipio. Essas extens\~oes podem fortalecer a robustez metodol\'ogica e a aplicabilidade do sistema em outros cen\'arios amaz\^onicos.
